\begin{theorem}[The positive-rung Technical Theorem in large order]
Let \(G=(V,E)\) be a finite loopless multigraph with its individual edge
identities, let \(e_0\in E\), and let the palette have \(k\) colors.  Let

\[
 \mathcal E:\noderef{ExtendedTashkinovTree}(G,k,e_0,\rho,T,\varphi)
\]

be a positive-rung ETT with displayed tree \(T\), generating coloring
\(\varphi\), rung number \(n\), and certified series data.  Assume
\noderef{MaximumProperty} for \(\mathcal E\) and

\[
 \noderef{MaximumDegree}(G)+1\leq k,
 \qquad 11\leq\noderef{TashkinovOrder}(G).
\]

For \(i\in\operatorname{Fin}(n+1)\), put
\[
 L_i=\mathcal E.\mathrm{series.data.tree}(i.\mathrm{castSucc}),\quad
 D_i=\noderef{TashkinovStoredColors}
        (\mathcal E.\mathrm{series.data},i),\quad
 \varphi_i=\mathcal E.\mathrm{series.data.coloring}(i).
\]
Write \(a=\operatorname{Fin.last}(n)\), and put
\[
 F=\mathcal E.\mathrm{series.data.tree}
       (\operatorname{Fin.last}(n+1)).
\]
Thus \(L_a\) is the exact terminal ladder tree and \(F\) is the exact final
series tree.  Write \(\Theta_i,u_i,\gamma_i,\delta_i\) for the kind, chosen
vertex, gamma color, and delta color stored at index \(i\), and write
\(S_i^{\mathrm{raw}}\) for the raw stored set at \(i\).  Then all of the
following assertions hold.

\begin{enumerate}
\item[(i)]
\(\noderef{ColoredVertexSetElementary}
  (G,\varphi,\noderef{TreeSequenceVertexSet}(T))\).

\item[(ii)]
For every partial coloring \(\sigma\), if \(\sigma\) is
\noderef{IsProperEdgeColoring}, leaves exactly \(e_0\) uncolored in the
sense of \noderef{IsOneEdgeUncolored}, and is
\noderef{StableColoring} on \(F\) relative to \((D_a,\varphi_a)\), then,
for all colors \(\alpha,\beta\),
\[
 \alpha\in\noderef{TreeMissingColors}(\sigma,F)
 \quad\hbox{or}\quad
 \beta\in\noderef{TreeMissingColors}(\sigma,F)
 \quad\Longrightarrow\quad
 \noderef{TreeInterchangeable}(G,F,\sigma,\alpha,\beta).
\]

\item[(iii)]
For every \(q\in\operatorname{Fin}(n)\), vertex \(v\), certified witness
\(h:v\in\noderef{TreeSequenceVertexSet}(T)\), and
\(j\in\operatorname{Fin}(n+1)\), suppose
\[
 \Theta_{q^+}=\mathrm{PE},\qquad u_{q^+}=v,\qquad
 \noderef{ETTFirstAppearanceIndex}(\mathcal E,v,h)=j+1.
\]
Every proper coloring \(\sigma\) leaving exactly \(e_0\) uncolored which is
stable on \(L_{q^+}\) relative to \((D_{q^+},\varphi_{q^+})\) is
\noderef{StableColoringBeforeVertex} before \(v\) in \(T\), using the same
witness \(h\), relative to \((D_j,\varphi_j)\).  Moreover, if \(j>0\),
then \(\sigma\) is stable on \(L_j\) relative to
\((D_j,\varphi_j)\).

In addition, let \(q,r\in\operatorname{Fin}(n)\) be distinct, and let
\(v_q,v_r\) be distinct vertices with certified membership witnesses in
\(T\).  If
\[
 \Theta_{q^+}=\Theta_{r^+}=\mathrm{PE},\qquad
 u_{q^+}=v_q,\qquad u_{r^+}=v_r,
\]
and their two dependent
\noderef{ETTFirstAppearanceIndex} values are equal, then
\[
          S_{q^+}^{\mathrm{raw}}\cap S_{r^+}^{\mathrm{raw}}
          =\varnothing.
\]

\item[(iv)]
For all \(v\in V\) and colors \(\gamma,\delta\), suppose
\[
 \Theta_a=\mathrm{PE},\qquad u_a=v,\qquad
 \gamma_a=\gamma,\qquad\delta_a=\delta.
\]
Under every proper, exactly-\(e_0\)-uncolored coloring \(\sigma\) stable on
\(L_a\) relative to \((D_a,\varphi_a)\), every \(w\in V(L_a)\) satisfies
\[
 \operatorname{inl}(w)\in
   \noderef{TwoColorComponent}
     (G,\sigma,\gamma,\delta,\operatorname{inl}(v))
 \quad\Longleftrightarrow\quad w=v.
\]

\item[(v)]
Let \(\sigma\) be any proper, exactly-\(e_0\)-uncolored coloring stable on
\(L_a\) relative to \((D_a,\varphi_a)\), and let \(d\) be a
\noderef{DefectiveColor} of \(V(L_a)\) under \(\sigma\).  Let \(v\) satisfy
\noderef{BoundaryColorVertex}\((G,\sigma,V(L_a),d,v)\), and suppose there
is another such boundary-color vertex \(w\) with
\noderef{TreeVertexOrder}\((L_a,w,v)\).  For every certified witness
\(h:v\in V(T)\), every \(q\in\operatorname{Fin}(n)\), and every vertex
\(v_i\), if
\[
 \Theta_{q^+}=\mathrm{PE}\ \hbox{or}\ \Theta_{q^+}=\mathrm{SE},
 \qquad u_{q^+}=v_i,\qquad
 \noderef{ETTFirstAppearanceIndex}(\mathcal E,v,h)\leq q+1,
\]
then
\[
                v=v_i\quad\hbox{or}\quad
                \noderef{TreeVertexOrder}(T,v,v_i).
\]

\item[(vi)]
Every proper, exactly-\(e_0\)-uncolored coloring \(\sigma\) stable on
\(L_a\) relative to \((D_a,\varphi_a)\) satisfies
\(\noderef{ModColoring}(\mathcal E,\sigma)\).
\end{enumerate}
\end{theorem}
\begin{proof}
We first construct, inside this proof, a
\noderef{TechnicalInductionContext} \(C_b\) for every natural bound \(b\).
We use strong induction on \(b\).  Fix \(b\), and suppose that \(C_s\) has
already been constructed for every \(s<b\).  To form \(C_b\), begin by
establishing all six assertions for an arbitrary positive-rung maximum ETT
\(\mathcal A\) whose rung number is \(s<b\), allowing its certified input
coloring to be arbitrary.

The strong-induction hypothesis gives \(C_s\).  Applying
\noderef{TechnicalCurrentRungElementaryLarge} to \(\mathcal A\), its
maximum property, \(C_s\), and the two numerical assumptions proves
assertion (i) at rung \(s\).

For assertion (ii), apply
\noderef{RutgersCurrentRungInterchangeability} to \(\mathcal A\).  Its
four induction premises follow from \(C_s\) as follows.  For any
arbitrary-input maximum ETT whose rung equals \(s\), transport \(C_s\)
along that rung equality and apply
\noderef{TechnicalCurrentRungElementaryLarge}; this gives current
assertion (i) uniformly.  The fields \texttt{currentFour} and
\texttt{currentSix} of \(C_s\) give current assertions (iv) and (vi) for
the same arbitrary-input ETTs.  Finally, \texttt{lowerTwo} gives assertion
(ii) for every positive maximum ETT of rung strictly below \(s\).  These
are exactly the current-(i),(iv),(vi) and strict-lower-(ii) premises of the
Rutgers lemma, so its conclusion is assertion (ii) for \(\mathcal A\),
with the exact final series tree and terminal stability data.

The remaining four assertions for \(\mathcal A\) are the fields
\texttt{currentThree}, \texttt{currentFour}, \texttt{currentFive}, and
\texttt{currentSix} of \(C_s\).  We have therefore proved assertions
(i)--(vi) for every positive-rung maximum ETT of rung \(s<b\).  These six
results define the fields \texttt{lowerOne} through \texttt{lowerSix} of
\(C_b\), preserving all dependent indices, membership witnesses, exact
trees, stored sets, and colorings appearing in the definition of the
context.

It remains to define the four current-rung fields of \(C_b\).  Fix an
arbitrary-input maximum ETT of rung exactly \(b\).  The complete
strict-lower package just proved supplies all six assertions for each
positive-rung maximum ETT below \(b\).  Apply
\noderef{ExtensionRules} to obtain current assertion (iii), including
both the stability conclusions and the disjointness of raw stored sets.
Apply \noderef{ExtensionBase} to obtain current assertion (iv), with the
exact terminal PE data.  Apply \noderef{Basic2} to obtain current
assertion (v), retaining the defective boundary identities, chosen
vertex, and first-appearance bound.  Apply \noderef{Td} to obtain current
assertion (vi), universally over the eligible terminal-stable colorings.
This constructs \(C_b\) and completes the strong induction.

Now take \(b=n+1\).  The positive-rung field of
\noderef{ExtendedTashkinovTree} gives \(0<n\), while arithmetic gives
\(n<n+1\).  Hence the given \(\mathcal E\), together with its maximum
property, satisfies the common premises of all six lower fields of
\(C_{n+1}\).  Field \texttt{lowerOne} is exactly (i), and
\texttt{lowerTwo} is exactly (ii): its terminal index is
\(a=\operatorname{Fin.last}(n)\), while its final-tree index is
\(\operatorname{Fin.last}(n+1)\).

Likewise, \texttt{lowerThree} is exactly the two-part assertion (iii),
including its dependent first-appearance witnesses and the \(j>0\) guard;
\texttt{lowerFour} is (iv) with the exact data
\(\Theta_a,u_a,\gamma_a,\delta_a\); \texttt{lowerFive} is (v) with the
same boundary-color and order witnesses; and \texttt{lowerSix} is (vi)
on \(L_a\) relative to \((D_a,\varphi_a)\).  Their conjunction proves all
six assertions.
\end{proof}
